\subsubsection{覆盖内容、原理与求解步骤}

规划模型先按变量和约束结构分类：连续线性规划、非线性规划、二次规划、整数
规划、混合整数线性规划和 0--1 规划；再区分单目标与多目标、确定性与随机性。
线性或凸结构优先使用有最优性证书的确定性求解器。整数变量表示启停、选择、
顺序和分配时，必须写出 Big-M 的来源或改用指示约束，不能以取整替代整数规划。
多目标问题先报告 Pareto 关系，再依据题意选择加权、分层或 $\varepsilon$ 约束，
不凭个人偏好给权重。

智能优化覆盖遗传算法、粒子群、模拟退火、蚁群、差分进化和贪心基线，以及
NSGA-II、MOPSO。思维导图中的鲸鱼优化、麻雀搜索、蛾火焰优化作为受控对照：
只有在变量编码、可行性修复和公平预算明确后才进入实验。复杂名称不构成创新，
必须与确定性基线或简单启发式在同一目标、随机种子组和函数评估次数下比较。

图论部分覆盖 Dijkstra、Floyd、SPFA，Prim、Kruskal，最大流、最小费用流、
路径规划、二分图匹配和节点重要性。稀疏非负边单源问题优先 Dijkstra；全源且
规模中等用 Floyd；存在负边时先检查负环，SPFA 只作工程候选而非复杂度保证。
最小生成树解决连通成本，不替代最短路；最大流回答容量，最小费用流同时回答
流量和成本；指派关系优先二分图匹配。

统一求解步骤为：
\begin{enumerate}
  \item 定义决策变量、目标、硬约束、软约束和不确定参数；
  \item 建立可行性基线，并在小规模枚举样例上核对目标值；
  \item 判断线性、凸性、整数性、网络结构和多目标冲突；
  \item 先运行确定性/贪心基线，再决定是否启用元启发式；
  \item 写清编码、初始化、选择更新、约束修复、终止条件和随机种子；
  \item 记录可行率、最优界/间隙、收敛轨迹、运行时间和跨种子分布；
  \item 通过 Monte Carlo 或情景仿真传播参数不确定性，并做敏感性分析。
\end{enumerate}

